\documentclass[a4paper,10pt]{article}
\usepackage[T1]{fontenc}
\usepackage{lmodern}
\usepackage[margin=21mm,headheight=14pt]{geometry}
\usepackage{amsmath,amssymb,booktabs,array,microtype}
\usepackage{xcolor,fancyhdr,hyperref}
\definecolor{ink}{HTML}{183E52}
\hypersetup{colorlinks=true,urlcolor=ink,linkcolor=ink,
 pdftitle={Speziale 2001 MgO EOS: Debye-temperature calculation},pdfauthor={Clemens Prescher}}
\pagestyle{fancy}
\fancyhf{}
\fancyhead[L]{\small Speziale (2001) MgO EOS}
\fancyhead[R]{\small Peritheos / 9 October 2026}
\fancyfoot[C]{\small\thepage}
\renewcommand{\headrulewidth}{0.3pt}
\setlength{\parindent}{0pt}
\setlength{\parskip}{5pt}
\setlength{\abovedisplayskip}{6pt}
\setlength{\belowdisplayskip}{6pt}
\newcommand{\sectiontitle}[1]{\vspace{5pt}{\large\bfseries\color{ink}#1}\par}
\newcommand{\GPa}{\mathrm{GPa}}
\begin{document}
{\LARGE\bfseries\color{ink}Reconstructing the MgO thermal EOS}\par
{\large Equations and Debye-temperature prescription in Speziale (2001)}\par
\vspace{3pt}
\textbf{Purpose.} Reconstruct the thermal EOS step by step, starting from the
300 K reference curve. Sergio Speziale has confirmed the Debye-temperature
integration used for the final variable-$q$ model (see Section 7).
All equations needed for our calculation are
written out here. Labels S1--S4, S10 and S11 correspond directly to the same
equation numbers in Speziale et al. (2001), expressed in the notation below.
Labels N1--N11 identify equations or derivations added in this note.

\sectiontitle{1. Reference state and normalization}
For $\gamma_0$, we use the value explicitly fixed in Section 5.1 of the paper:
$\gamma_0=1.524$. The Figure 6 caption reports $\gamma_0=1.524\pm0.025$.
All calculations below use this printed central value; $\gamma_0$ is not refitted.
\begin{center}
\renewcommand{\arraystretch}{1.14}
\begin{tabular}{lll}
\toprule
Quantity & Adopted value & Definition \\
\midrule
$V_0$ & $74.71\ \text{\AA}^{3}$ & Conventional-cell volume at $T_r$, zero pressure \\
$K_0$ & $160.2\ \GPa$ & Isothermal bulk modulus at the reference state \\
$K'_0$ & $3.99$ & $(\partial K_T/\partial P)_{T_r}$ at zero pressure \\
$\gamma_0$ & $1.524$ & Published fixed value (Section 5.1) \\
$q_0,\ q_1$ & $1.65,\ 11.8$ & Constants in the variable-$q$ model \\
$\theta_0$ & $773\ \mathrm{K}$ & Reference Debye temperature \\
$T_r$ & $300\ \mathrm{K}$ & Principal isotherm and energy subtraction \\
$Z,\ n$ & $4,\ 2$ & Formula units per cell; atoms per MgO formula unit \\
\bottomrule
\end{tabular}
\end{center}
$V_c$ denotes the numerical conventional-cell volume in \AA$^3$;
$V_m$ is the molar volume in m$^3$ per mole of MgO. Define
\begin{equation}\label{eq:volume}
 x=\frac{V_c}{V_0},\qquad V_m=\frac{N_A}{Z}\,V_c\,10^{-30},
 \qquad N_A=6.02214076\times10^{23}\ \mathrm{mol}^{-1}.\tag{N1}
\end{equation}
This unit conversion is explicit here; it is not a numbered equation in the
paper. The gas constant is $R=8.31446261815324\ \mathrm{J\,mol^{-1}\,K^{-1}}$.

\sectiontitle{2. Complete pressure calculation}
Sections 2.1 and 4.2 specify a third-order Birch--Murnaghan reference isotherm.
We write its standard expression explicitly:
\begin{equation}\label{eq:bm3}
 P_{300}(V_c)=\frac{3K_0}{2}
 \left(x^{-7/3}-x^{-5/3}\right)
 \left[1+\frac{3}{4}(K'_0-4)\left(x^{-2/3}-1\right)\right].\tag{N2}
\end{equation}
With $K_0$ in GPa, this expression returns GPa. For the Debye energy, the
paper's Equation 3 becomes, with $y=\theta(V_c)/T$,
\begin{equation}\label{eq:energy}
 E_D(V_c,T)=\frac{9nRT}{y^3}\int_0^y\frac{z^3}{e^z-1}\,dz.\tag{S3}
\end{equation}
$E_D$ is in J per mole of MgO; $n=2$. The same $\theta(V_c)$ is used at both
temperatures in the fixed-volume subtraction. The paper's Equations 2 and 1
then give the thermal increment and total pressure, respectively:
\begin{align}
 \Delta P_{\mathrm{th}}(V_c,T)
 &=10^{-9}\frac{\gamma(V_c)}{V_m}
 \left[E_D(V_c,T)-E_D(V_c,T_r)\right],\tag{S2}\label{eq:thermal}\\
 P(V_c,T)&=P_{300}(V_c)+\Delta P_{\mathrm{th}}(V_c,T).\tag{S1}\label{eq:total}
\end{align}
The factor $10^{-9}$ converts Pa to GPa. Thus $P(V_0,300\ \mathrm{K})=0$.
The electronic contribution is omitted, as in the MgO model under discussion;
any temperature-independent vibrational term cancels in the energy difference.

\newpage
\sectiontitle{3. Volume dependence of the Gr\"uneisen parameter}
Section 2.1 defines $q$ through the logarithmic volume derivative of $\gamma$.
We write that unnumbered definition as
\begin{equation}\label{eq:qdef}
 q(V_c)=\left(\frac{\partial\ln\gamma}{\partial\ln V_c}\right)_T.\tag{N3}
\end{equation}
The paper's final variable-$q$ model (its Equations 10--11) is
\begin{align}
 q(V_c)&=q_0x^{q_1},\tag{S10}\label{eq:q}\\
 \gamma(V_c)&=\gamma_0\exp\!\left[
 \frac{q_0}{q_1}\left(x^{q_1}-1\right)\right].\tag{S11}\label{eq:gamma}
\end{align}
Equation~\eqref{eq:gamma} follows by integrating Equation~\eqref{eq:qdef} with
$\gamma(V_0)=\gamma_0$. It also gives a finite value at infinite compression:
\begin{equation}\label{eq:glimit}
 \lim_{x\to0}\gamma(x)=\gamma_0e^{-q_0/q_1}=1.325127\quad
 \text{for the adopted coefficients}.\tag{N4}
\end{equation}

\sectiontitle{4. Debye-temperature calculation and the question}
We reconstruct the Debye temperature from the Debye--Gr\"uneisen relation
below. This is the thermodynamic relation used for our derivation, rather
than an explicit numbered equation in the paper:
\begin{equation}\label{eq:thetadef}
 \gamma(V_c)=-\frac{d\ln\theta(V_c)}{d\ln V_c},
 \qquad\theta(V_0)=\theta_0.\tag{N5}
\end{equation}
For \textbf{constant} $q\ne0$, integration gives
\begin{align}
 \gamma(V_c)&=\gamma_0x^q,\tag{N6}\label{eq:gammaconstant}\\
 \theta(V_c)&=\theta_0\exp\!\left[\frac{\gamma_0-\gamma(V_c)}{q}\right].\tag{S4}\label{eq:constant}
\end{align}
The second expression is the paper's Equation 4; the first is the constant-$q$
gamma law discussed in Section 5.4. For $q=0$, the corresponding
limit is $\gamma(V_c)=\gamma_0$ and $\theta(V_c)=\theta_0x^{-\gamma_0}$.

For the \textbf{variable-$q$} model, our implementation instead integrates
Equation~\eqref{eq:thetadef} using Equation~\eqref{eq:gamma}:
\begin{equation}\label{eq:thetaint}
 \boxed{\displaystyle
 \theta(V_c)=\theta_0\exp\!\left[
 -\int_1^x\gamma_0
 \exp\!\left(\frac{q_0}{q_1}(u^{q_1}-1)\right)\frac{du}{u}
 \right].}\tag{N7}
\end{equation}
An independent closed form, with $a=q_0/q_1$, is
\begin{equation}\label{eq:ei}
 \ln\frac{\theta(V_c)}{\theta_0}
 =-\frac{\gamma_0e^{-a}}{q_1}
 \left[\operatorname{Ei}(ax^{q_1})-\operatorname{Ei}(a)\right].\tag{N8}
\end{equation}
Here $\operatorname{Ei}$ is the exponential integral, defined on positive
arguments by $\operatorname{Ei}(z)=\mathrm{PV}\int_{-\infty}^{z}e^t/t\,dt$,
where PV denotes the Cauchy principal value. Its derivative is
$d\operatorname{Ei}(z)/dz=e^z/z$. Equation~\eqref{eq:ei} is used for the
adopted positive $q_0$ and $q_1$.

\textbf{Diagnostic illustration, not a known author prescription.}
Simply inserting instantaneous $q(V_c)$ into Equation~\eqref{eq:constant}
would give the hypothetical expression
\begin{equation}\label{eq:localtheta}
 \theta_{\mathrm{local}}(V_c)
 =\theta_0\exp\!\left[\frac{\gamma_0-\gamma(V_c)}{q(V_c)}\right].\tag{N9}
\end{equation}
For Equations~\eqref{eq:q}--\eqref{eq:gamma}, differentiating this candidate gives
\begin{equation}\label{eq:localderivative}
 -\frac{d\ln\theta_{\mathrm{local}}}{d\ln V_c}
 =\gamma(V_c)+\frac{q_1}{q(V_c)}\left[\gamma_0-\gamma(V_c)\right].\tag{N10}
\end{equation}
The extra term arises because $q(V_c)$ also changes with volume. This
\emph{hypothetical expression}, Equation~\eqref{eq:localtheta}, therefore does
not satisfy Equation~\eqref{eq:thetadef} in general. Our integrated
Equation~\eqref{eq:thetaint} does satisfy that relation. Speziale's reply now
confirms that this integration was used in his original calculation.

\newpage
\sectiontitle{5. Reference calculations}
These \textbf{Peritheos calculations} use the equations on pages 1--2 and the
published central parameter values. They are model outputs, not measured or
transcribed observations. The digits facilitate numerical comparison and do
not indicate physical parameter accuracy.
\begin{center}
\renewcommand{\arraystretch}{1.3}
\setlength{\tabcolsep}{5pt}
\begin{tabular}{rrrrr}
\toprule
$V_c/V_0$ & $V_c$ (\AA$^3$) & $T$ (K) & $\gamma$ & $\theta$ (K) \\
\midrule
1.00 & 74.71000 & 300  & 1.524000 & 773.000 \\
1.00 & 74.71000 & 1100 & 1.524000 & 773.000 \\
0.90 & 67.23900 & 1100 & 1.379666 & 899.275 \\
0.80 & 59.76800 & 2000 & 1.338508 & 1054.804 \\
0.64 & 47.81440 & 3663 & 1.326084 & 1419.208 \\
\bottomrule
\end{tabular}

\vspace{8pt}
\begin{tabular}{rrrrr}
\toprule
$V_c/V_0$ & $T$ (K) & $P_{300}$ (GPa) & $\Delta P_{\mathrm{th}}$ (GPa) & $P$ (GPa) \\
\midrule
1.00 & 300  & 0.000000   & 0.000000  & 0.000000 \\
1.00 & 1100 & 0.000000   & 4.964240  & 4.964240 \\
0.90 & 1100 & 20.830878  & 4.855028  & 25.685906 \\
0.80 & 2000 & 55.839883  & 11.610349 & 67.450232 \\
0.64 & 3663 & 174.738015 & 28.659685 & 203.397700 \\
\bottomrule
\end{tabular}
\end{center}

\sectiontitle{6. Checks and their limits}
\textbf{Equation and normalization checks.} We verify the defining derivatives
in Equations~\eqref{eq:qdef} and \eqref{eq:thetadef}. At $T=T_r$, the thermal
pressure is zero at every volume, and $P(V_0,T_r)=0$. Evaluating the Debye
energy per mole of MgO ($n=2$, $Z=4$) or per mole of conventional cells
($n=8$, $Z=1$) gives the same pressure. The high-temperature limit is
\begin{equation}\label{eq:highT}
 \lim_{T/\theta\to\infty}\frac{E_D(V_c,T)}{T}=3nR=6R
 \quad\text{per mole of MgO}.\tag{N11}
\end{equation}

\textbf{Independent numerical evaluation.} With all coefficients fixed to the
published central values, the 64-point Gauss--Legendre calculation, independently
coded adaptive integration and the closed form in Equation~\eqref{eq:ei} agree
to better than $10^{-9}$ GPa on 35 states spanning $x=0.60$--$1.02$ and
$T=300$--$3663$ K. This checks numerical evaluation of the stated equations;
it is not a claim of physical accuracy at that pressure precision.

\textbf{Comparison with the paper's calculated curves.} Figure 6 tests the
constant-$q$ branch ($q=1.65$), using Equation~\eqref{eq:constant}. Eleven
selected checkpoints on the 300, 1100 and 3000 K curves differ by at most
approximately 0.11 GPa, or 0.0004 in $V_c/V_0$. The Figure 10 inset tests the
final variable-$q$ branch: four selected checkpoints on its 1100 K curve differ
by at most approximately 0.30 GPa, or 0.0012 in $V_c/V_0$. These are comparisons
with digitized \emph{calculated curves}, not refits to measurements. Scan
resolution, axis placement, dotted-line phase and stroke width limit this
comparison. The integration prescription is now author-confirmed; these
graphical differences do not recover the original numerical settings.

\textbf{Remaining scope.} The complete original thermal fitting procedure and
shock-data reduction have not been reproduced. The integration prescription
is confirmed, but the original calculator, trapezoidal grid and numerical
settings remain unavailable.

\sectiontitle{7. Author confirmation recorded on 9 October 2026}
In the reply supplied by Clemens, Sergio Speziale states:
\emph{``I can confirm that I numerically integrated your equation N5 by using
S11, using the equation numbering of your notes.''} He adds:
\emph{``I generally used a simple trapezoidal integration approach.''}
He cannot currently retrieve the original calculator. This confirms our
continuous integration prescription, while leaving exact historical numerical
replay and the complete author fit unestablished.

\vspace{4pt}
{\small\textbf{Primary reference.} Speziale, Zha, Duffy, Hemley and Mao (2001),
\emph{Quasi-hydrostatic compression of magnesium oxide to 52 GPa:
Implications for the pressure-volume-temperature equation of state},
J. Geophys. Res. 106, 515-528.
\href{https://doi.org/10.1029/2000JB900318}{doi:10.1029/2000JB900318}.
Source model: Section 2.1, Equations 1-4; Sections 4.2 and 5.1-5.4,
Equations 10-11; Figures 6, 9 and 10.}
\end{document}
